% GENERATED FILE. Edit canonical lessons, then run npm run generate:books.
% canonical-source-hash: fnv1a64:cff245b97c59d711
% canonical-lessons: JA-C357-keiyaku-suru, JA-C357-hoshou, JA-C357-naiyou, JA-C357-yutori, JA-C357-kettei-suru

\chapter{Contracts}
\label{ch:ja-contracts}
\hlchaptermodality{357}

\begin{chapteropening}
\textbf{By the end of this chapter:} \emph{I can say when I will sign a contract and what it covers.}

Complete the last lesson of chapter 357: I can say when I will sign a contract and what it covers.
\end{chapteropening}

\section[\texorpdfstring{\ja{けいやくする} (keiyaku suru)}{けいやくする (keiyaku suru)}]{\ja{けいやくする} (keiyaku suru) — to sign a contract}

\label{lesson:JA-C357-keiyaku-suru}



\begin{quote}
Before the new one: say the Japanese for in most cases, then the Japanese for forever.
\end{quote}

\subsection*{You'll want to know: \ja{けいやくする}}
\textbf{\ja{けいやくする}} — \emph{keiyaku suru} — \textquotedblleft{}to sign a contract\textquotedblright{}.

\begin{quote}
\textbf{\ja{じょうけんに} \ja{さんせいしたら、けいやくします。}} — \emph{jōken ni sansei shitara, keiyaku shimasu} — \textquotedblleft{}If I agree to the terms, I will sign the contract.\textquotedblright{}
\end{quote}

Polite \textbf{\ja{けいやくします}}, \emph{keiyakushimasu}; a \ja{する} verb. Past (polite): \textbf{\ja{ちちが} \ja{かいしゃと} \ja{けいやくしました。}}, \emph{chichi ga kaisha to keiyaku shimashita}.

\begin{cousinweb}
\textbf{\ja{けいやくする}} (\emph{keiyaku suru}) is built from a Sino-Japanese noun and the native verb \emph{suru}, \textquotedblleft{}to do\textquotedblright{}. The noun, \emph{keiyaku}, is two Chinese characters read with their on'yomi, readings borrowed from Chinese: \emph{kei}, \textquotedblleft{}pledge\textquotedblright{} (Middle Chinese \emph{khejH}) and \emph{yaku}, \textquotedblleft{}promise\textquotedblright{} (\emph{'jak}) (KANJIDIC2; JMdict; Middle Chinese in Baxter's transcription).
\end{cousinweb}

\begin{grammarlens}[title={What you've built}]
One more word for saying what happens if.
\end{grammarlens}

\subsection*{Your turn}
\begin{itemize}
\raggedright
  \item \emph{Say it:} \emph{keiyaku suru}
  \item \emph{Say it:} \emph{keiyaku suru}, once more
  \item \emph{Say it:} \emph{jōken ni sansei shitara, keiyaku shimasu} — If I agree to the terms, I will sign the contract.
  \item \emph{Recall:} say the Japanese for in most cases, then the Japanese for forever, then say \emph{keiyaku suru} again
\end{itemize}

\begin{tcolorbox}[breakable,colback=teal!4,colframe=teal!35!black,title={Before you move on}]
What does \ja{けいやくする} mean? (\textquotedblleft{}To sign a contract\textquotedblright{}.) Say it once more.
\end{tcolorbox}
\section[\texorpdfstring{\ja{ほしょう} (hoshō)}{ほしょう (hoshō)}]{\ja{ほしょう} (hoshō) — a guarantee; a warranty}

\label{lesson:JA-C357-hoshou}



\begin{quote}
Before the new one: say the Japanese for forever, then the Japanese for to sign a contract.
\end{quote}

\subsection*{You'll want to know: \ja{ほしょう}}
\textbf{\ja{ほしょう}} — \emph{hoshō} — \textquotedblleft{}a guarantee; a warranty\textquotedblright{}.

\begin{quote}
\textbf{\ja{こわれたら、ほしょうが} \ja{ありますか。}} — \emph{kowaretara, hoshō ga arimasu ka} — \textquotedblleft{}If it breaks, is there a warranty?\textquotedblright{}
\end{quote}

\begin{cousinweb}
\textbf{\ja{ほしょう}} (\emph{hoshō}) is Sino-Japanese (kango): two Chinese characters read with their on'yomi, readings borrowed from Chinese. They are \emph{ho}, \textquotedblleft{}preserve\textquotedblright{} (Middle Chinese \emph{pawX}) and \emph{shō}, \textquotedblleft{}evidence\textquotedblright{} (\emph{tsyingH}) (KANJIDIC2; Middle Chinese in Baxter's transcription).
\end{cousinweb}

\begin{grammarlens}[title={What you've built}]
One more word for saying what happens if.
\end{grammarlens}

\subsection*{Your turn}
\begin{itemize}
\raggedright
  \item \emph{Say it:} \emph{hoshō}
  \item \emph{Say it:} \emph{hoshō}, once more
  \item \emph{Say it:} \emph{kowaretara, hoshō ga arimasu ka} — If it breaks, is there a warranty?
  \item \emph{Recall:} say the Japanese for forever, then the Japanese for to sign a contract, then say \emph{hoshō} again
\end{itemize}

\begin{tcolorbox}[breakable,colback=teal!4,colframe=teal!35!black,title={Before you move on}]
What does \ja{ほしょう} mean? (\textquotedblleft{}A guarantee; a warranty\textquotedblright{}.) Say it once more.
\end{tcolorbox}
\section[\texorpdfstring{\ja{ないよう} (naiyō)}{ないよう (naiyō)}]{\ja{ないよう} (naiyō) — the contents; the details}

\label{lesson:JA-C357-naiyou}



\begin{quote}
Before the new one: say the Japanese for to sign a contract, then the Japanese for a warranty.
\end{quote}

\subsection*{You'll want to know: \ja{ないよう}}
\textbf{\ja{ないよう}} — \emph{naiyō} — \textquotedblleft{}the contents; the details\textquotedblright{}.

\begin{quote}
\textbf{\ja{ないようを} \ja{よんだら、はんこを} \ja{おして} \ja{ください。}} — \emph{naiyō o yondara, hanko o oshite kudasai} — \textquotedblleft{}Once you have read the contents, please stamp it with your seal.\textquotedblright{}
\end{quote}

\begin{cousinweb}
\textbf{\ja{ないよう}} (\emph{naiyō}) is Sino-Japanese (kango): two Chinese characters read with their on'yomi, readings borrowed from Chinese. They are \emph{nai}, \textquotedblleft{}inside\textquotedblright{} (Middle Chinese \emph{nwojH}) and \emph{yō}, \textquotedblleft{}contain\textquotedblright{} (\emph{yowng}) (KANJIDIC2; Middle Chinese in Baxter's transcription).
\end{cousinweb}

\begin{grammarlens}[title={What you've built}]
One more word for saying what happens if.
\end{grammarlens}

\subsection*{Your turn}
\begin{itemize}
\raggedright
  \item \emph{Say it:} \emph{naiyō}
  \item \emph{Say it:} \emph{naiyō}, once more
  \item \emph{Say it:} \emph{naiyō o yondara, hanko o oshite kudasai} — Once you have read the contents, please stamp it with your seal.
  \item \emph{Recall:} say the Japanese for to sign a contract, then the Japanese for a warranty, then say \emph{naiyō} again
\end{itemize}

\begin{tcolorbox}[breakable,colback=teal!4,colframe=teal!35!black,title={Before you move on}]
What does \ja{ないよう} mean? (\textquotedblleft{}The contents; the details\textquotedblright{}.) Say it once more.
\end{tcolorbox}
\section[\texorpdfstring{\ja{ゆとり} (yutori)}{ゆとり (yutori)}]{\ja{ゆとり} (yutori) — leeway; room to spare}

\label{lesson:JA-C357-yutori}



\begin{quote}
Before the new one: say the Japanese for a warranty, then the Japanese for the contents.
\end{quote}

\subsection*{You'll want to know: \ja{ゆとり}}
\textbf{\ja{ゆとり}} — \emph{yutori} — \textquotedblleft{}leeway; room to spare\textquotedblright{}.

\begin{quote}
\textbf{\ja{おかねに} \ja{ゆとりが} \ja{あったら、ちょきんします。}} — \emph{okane ni yutori ga attara, chokin shimasu} — \textquotedblleft{}If I have money to spare, I save it.\textquotedblright{}
\end{quote}

\begin{cousinweb}
\textbf{\ja{ゆとり}} (\emph{yutori}) is a native Japanese word (wago), inherited rather than borrowed from Chinese, and its older origin is not settled; JMdict gives it no kanji spelling (JMdict).
\end{cousinweb}

\begin{grammarlens}[title={What you've built}]
One more word for saying what happens if.
\end{grammarlens}

\subsection*{Your turn}
\begin{itemize}
\raggedright
  \item \emph{Say it:} \emph{yutori}
  \item \emph{Say it:} \emph{yutori}, once more
  \item \emph{Say it:} \emph{okane ni yutori ga attara, chokin shimasu} — If I have money to spare, I save it.
  \item \emph{Recall:} say the Japanese for a warranty, then the Japanese for the contents, then say \emph{yutori} again
\end{itemize}

\begin{tcolorbox}[breakable,colback=teal!4,colframe=teal!35!black,title={Before you move on}]
What does \ja{ゆとり} mean? (\textquotedblleft{}Leeway; room to spare\textquotedblright{}.) Say it once more.
\end{tcolorbox}
\section[\texorpdfstring{\ja{けっていする} (kettei suru)}{けっていする (kettei suru)}]{\ja{けっていする} (kettei suru) — to decide officially}

\label{lesson:JA-C357-kettei-suru}



\begin{quote}
Before the new one: say the Japanese for the contents, then the Japanese for leeway.
\end{quote}

\subsection*{You'll want to know: \ja{けっていする}}
\textbf{\ja{けっていする}} — \emph{kettei suru} — \textquotedblleft{}to decide officially\textquotedblright{}.

\begin{quote}
\textbf{\ja{みんなが} \ja{さんせいしたら、けっていします。}} — \emph{minna ga sansei shitara, kettei shimasu} — \textquotedblleft{}If everyone agrees, it will be decided.\textquotedblright{}
\end{quote}

Polite \textbf{\ja{けっていします}}, \emph{ketteishimasu}; a \ja{する} verb. Past (polite): \textbf{\ja{かいしゃが} \ja{けっていしました。}}, \emph{kaisha ga kettei shimashita}.

\begin{cousinweb}
\textbf{\ja{けっていする}} (\emph{kettei suru}) is built from a Sino-Japanese noun and the native verb \emph{suru}, \textquotedblleft{}to do\textquotedblright{}. The noun, \emph{kettei}, is two Chinese characters read with their on'yomi, readings borrowed from Chinese: \emph{ketsu}, \textquotedblleft{}decide\textquotedblright{} (Middle Chinese \emph{kwet}) and \emph{tei}, \textquotedblleft{}determine\textquotedblright{} (\emph{dengH}) (KANJIDIC2; JMdict; Middle Chinese in Baxter's transcription). The two join as \emph{ket-tei}.
\end{cousinweb}

\begin{grammarlens}[title={What you've built}]
One more word for saying what happens if.
\end{grammarlens}

\subsection*{Your turn}
\begin{itemize}
\raggedright
  \item \emph{Say it:} \emph{kettei suru}
  \item \emph{Say it:} \emph{kettei suru}, once more
  \item \emph{Say it:} \emph{minna ga sansei shitara, kettei shimasu} — If everyone agrees, it will be decided.
  \item \emph{Recall:} say the Japanese for the contents, then the Japanese for leeway, then say \emph{kettei suru} again
\end{itemize}

\begin{tcolorbox}[breakable,colback=teal!4,colframe=teal!35!black,title={Before you move on}]
What does \ja{けっていする} mean? (\textquotedblleft{}To decide officially\textquotedblright{}.) Say it once more.
\end{tcolorbox}
